
\documentclass{article}
\usepackage{colortbl}
\usepackage{makecell}
\usepackage{multirow}
\usepackage{supertabular}

\begin{document}

\newcounter{utterance}

\centering \large Interaction Transcript for game `taboo', experiment `high\_ru', episode 15 with Qwen3.6{-}35B{-}A3B{-}FP8{-}without{-}reasoning.
\vspace{24pt}

{ \footnotesize  \setcounter{utterance}{1}
\setlength{\tabcolsep}{0pt}
\begin{supertabular}{c@{$\;$}|p{.15\linewidth}@{}p{.15\linewidth}p{.15\linewidth}p{.15\linewidth}p{.15\linewidth}p{.15\linewidth}}
    \# & $\;$A & \multicolumn{4}{c}{Game Master} & $\;\:$B\\
    \hline

    \theutterance \stepcounter{utterance}  
    & & \multicolumn{4}{p{0.6\linewidth}}{
        \cellcolor[rgb]{0.9,0.9,0.9}{
            \makecell[{{p{\linewidth}}}]{
                \texttt{\tiny{[P1$\langle$GM]}}
                \texttt{Вы играете в совместную игру на угадывание слов, в которой вам нужно описать целевое слово так, чтобы другой игрок смог его угадать.} \\
\\ 
\texttt{Правила:} \\
\texttt{(a) Вы должны отвечать в формате: ПОДСКАЗКА: <какой{-}то текст>. Попытки другого игрока будут начинаться со слова ДОГАДКА.} \\
\texttt{(b) Вы не можете использовать само целевое слово, его части или его морфологические варианты в своём описании.} \\
\texttt{(c) Кроме того, те же правила действуют для связанных слов, которые приведены ниже.} \\
\\ 
\texttt{Условия завершения:} \\
\texttt{(i) Если вы используете целевое слово или связанное слово в своём описании, то вы проигрываете.} \\
\texttt{(ii) Если другой игрок сможет угадать целевое слово за 3 попыток, то вы оба побеждаете.} \\
\\ 
\texttt{Давайте начнём.} \\
\\ 
\texttt{Вот целевое слово, которое вам нужно описать и которое другой игрок должен угадать:} \\
\\ 
\texttt{проект} \\
\\ 
\texttt{Связанные слова:} \\
\\ 
\texttt{{-} инициатива} \\
\texttt{{-} план} \\
\texttt{{-} намерение} \\
\\ 
\texttt{Важно: вы ограничены во времени, давайте короткие описания по существу!} \\
            }
        }
    }
    & \\ \\

    \theutterance \stepcounter{utterance}  
    & \multicolumn{4}{p{0.6\linewidth}}{
        \cellcolor[rgb]{0.9,0.9,0.9}{
            \makecell[{{p{\linewidth}}}]{
                \texttt{\tiny{[P1$\rangle$GM]}}
                \texttt{ПОДСКАЗКА: Это набор взаимосвязанных действий, направленных на достижение конкретной цели в рамках ограниченных ресурсов и сроков.} \\
            }
        }
    }
    & & \\ \\

    \theutterance \stepcounter{utterance}  
    & & & \multicolumn{2}{p{0.3\linewidth}}{
        \cellcolor[rgb]{0.9,0.9,0.9}{
            \makecell[{{p{\linewidth}}}]{
                \texttt{\tiny{[GM$|$GM]}}
                \texttt{continue} \\
            }
        }
    }
    & & \\ \\

    \theutterance \stepcounter{utterance}  
    & & & \multicolumn{2}{p{0.3\linewidth}}{
        \cellcolor[rgb]{0.9,0.9,0.9}{
            \makecell[{{p{\linewidth}}}]{
                \texttt{\tiny{[GM$|$GM]}}
                \texttt{это набор взаимосвязанных действий направленных на достижение конкретной цели в рамках ограниченных ресурсов и сроков} \\
            }
        }
    }
    & & \\ \\

    \theutterance \stepcounter{utterance}  
    & & \multicolumn{4}{p{0.6\linewidth}}{
        \cellcolor[rgb]{0.9,0.9,0.9}{
            \makecell[{{p{\linewidth}}}]{
                \texttt{\tiny{[GM$\rangle$P2]}}
                \texttt{Вы играете в совместную игру по угадыванию слов, в которой вам нужно угадать целевое слово, которое описывает вам другой игрок.} \\
\\ 
\texttt{В каждой попытке вы можете сделать одну догадку. Вы выигрываете, когда угадываете целевое слово. Вы проигрываете, если не можете угадать его за 3 попыток.} \\
\\ 
\texttt{После каждой попытки вы будете получать от другого игрока новую подсказку, которая начинается с ПОДСКАЗКА.} \\
\\ 
\texttt{Делайте свои догадки, просто называя слово в следующем формате: ДОГАДКА: <слово>} \\
\\ 
\texttt{Начнём.} \\
\\ 
\\ 
\texttt{ПОДСКАЗКА: Это набор взаимосвязанных действий, направленных на достижение конкретной цели в рамках ограниченных ресурсов и сроков.} \\
            }
        }
    }
    & \\ \\

    \theutterance \stepcounter{utterance}  
    & & & \multicolumn{4}{p{0.6\linewidth}}{
        \cellcolor[rgb]{0.9,0.9,0.9}{
            \makecell[{{p{\linewidth}}}]{
                \texttt{\tiny{[GM$\langle$P2]}}
                \texttt{ДОГАДКА: проект} \\
            }
        }
    }
     \\ \\

    \theutterance \stepcounter{utterance}  
    & & & \multicolumn{2}{p{0.3\linewidth}}{
        \cellcolor[rgb]{0.9,0.9,0.9}{
            \makecell[{{p{\linewidth}}}]{
                \texttt{\tiny{[GM$|$GM]}}
                \texttt{continue} \\
            }
        }
    }
    & & \\ \\

    \theutterance \stepcounter{utterance}  
    & & & \multicolumn{2}{p{0.3\linewidth}}{
        \cellcolor[rgb]{0.9,0.9,0.9}{
            \makecell[{{p{\linewidth}}}]{
                \texttt{\tiny{[GM$|$GM]}}
                \texttt{проект} \\
            }
        }
    }
    & & \\ \\

    \theutterance \stepcounter{utterance}  
    & & & \multicolumn{2}{p{0.3\linewidth}}{
        \cellcolor[rgb]{0.9,0.9,0.9}{
            \makecell[{{p{\linewidth}}}]{
                \texttt{\tiny{[GM$|$GM]}}
                \texttt{end game} \\
            }
        }
    }
    & & \\ \\

\end{supertabular}
}

\end{document}
